% Abhakseite „Das kann ich“ – Zone nur im Gesamt (\mitzone)
%% TODO Ich-kann-Sätze: Platzhalter wie in den Aufgabendateien
\begin{abhakseite}
\ifdefined\mitzone
\abhakgruppe{Kennst du schon}
\abhak{1}{Potenzen mit geradem und ungeradem Exponenten bei negativer Basis}
\abhak{2}{Öffnung einer Parabel und Verlauf der Normalparabel dritten Grades nach oben und unten}
\abhak{3}{Exponentialfunktion y = a · qˣ}
\abhak{4}{Achsen- und Punktsymmetrie am Graphen erkennen und f(−x) bilden (Vorzeichen bei geraden und ungeraden Potenzen)}
\abhak{5}{Satz vom Nullprodukt und die Regel, dass $e^x$ nie null und immer positiv ist}
\abhak{6}{Verschiebung und Streckung eines Graphen an der Gleichung lesen (f(x) + c verschiebt nach oben, a · f(x) streckt, negatives a spiegelt) und Monotonie als „steigt“ oder „fällt“}
\abhak{7}{Vorzeichen eines Produkts und eines Bruchs aus den Vorzeichen der Faktoren bestimmen (plus mal minus, Zähler positiv durch Nenner positiv)}
\abhak{8}{Vorzeichen eines Produkts und eines Bruchs aus den Vorzeichen der Faktoren bestimmen (plus mal minus, Zähler positiv durch Nenner positiv) – Fehler finden}
\abhak{9}{Vorzeichen eines Produkts und eines Bruchs aus den Vorzeichen der Faktoren bestimmen (plus mal minus, Zähler positiv durch Nenner positiv) – selbst rechnen}
\fi
\abhakgruppe{Einheit 1 · Ganzrationale Funktionen}
\abhak{10}{„Wer gewinnt?“}
\abhak{11}{Ganzrationale Funktionen}
\abhak{12}{Ganzrationale Funktionen – Fehler finden, Begründen, Darstellungswechsel}
\abhakgruppe{Einheit 2 · Produkte aus Polynom und e-Funktion}
\abhak{13}{Produkte aus Polynom und e-Funktion}
\abhak{14}{Produkte aus Polynom und e-Funktion – Fehler finden, Begründen, Anwendung, Darstellungswechsel}
\abhakgruppe{Einheit 3 · Waagerechte Asymptoten}
\abhak{15}{Waagerechte Asymptoten}
\abhak{16}{Waagerechte Asymptoten – Fehler finden, Begründen, Anwendung}
\end{abhakseite}
